% !TeX spellcheck = cs_CZ
\begin{mdframed}[style=mdexam]
  \begin{example}\label{mai:exam058}
    \textbf{Jak nekoupit zmetek}\newline
      Do finále soutěže o „šmejd roku“ se dostaly dva podniky, „Hvizd, s.r.o.“ a „Svist, a.s.“,
      které zásobují trh zábavnou pyrotechnikou. První z nich kryje požadavky trhu ze
      \qty{70}{\percent}, druhý ze zbývajících \qty{30}{\percent}. “Zjistilo se, že \qty{83}{\percent}
      ze všech výrobků Hvizdu je vadných (nebouchají, když je to třeba, zejména však bouchají, když
      se to nejméně očekává). V případě Svistu je zmetků pouze \qty{63}{\percent}. Porota soutěže
      rozhodla, že cenu dostane ten z obou podniků, jehož ředitel zodpoví správně následující
      otázky:

      {\centering
      \captionsetup{type=figure} 
      \includegraphics[width=0.8\linewidth]{mai_fig075.png}
      \captionof{figure}{Pat \& Mat: Silvestr}
      \label{mai_fig075}
      \par}

      \begin{enumerate}[noitemsep, label=\emph{\alph*}]
        \item Jaká je pravděpodobnost, že náhodně zakoupená rachejtle bude fungovat tak, jak má?
        \item Jaká je pravděpodobnost, že náhodně zakoupená rachejtle, o níž se na obalu píše, že
              byla vyrobena podnikem Hvizd, není zmetek?
        \item Jaká je pravděpodobnost, že náhodně zakoupená rachejtle, kterou se podařilo úspěšně
              odpálit, byla vyrobena podnikem Svist?
      \end{enumerate} 
      Postupně jednotlivé úkoly vyřešíme. V případě a) posuzujeme pravděpodobnost jevu \(A\):
      „Náhodně zakoupený výrobek je funkční.“ bez dalších podmínek. Jedná se o nepodmíněnou
      pravděpodobnost. Dejme tomu, že na trhu je v dané chvíli ke koupi \(n\) rachejtlí. Z nich
      \qty{70}{\percent}, tj. \(\num{0.7}n\), bylo vyrobeno v Hvizdu a zbytek, \(\num{0.3}n\), ve
      Svistu. Víme, že \qty{17}{\percent} výrobků Hvizdu je funkčních, v případě Svistu je to
      \qty{37}{\percent}. Na trhu je tedy v tuto chvíli
      \begin{equation*}
        m = \underbrace{\num{0.17}\cdot\num{0.7}n}_\text{Hvizd} + 
            \underbrace{\num{0.37}\cdot\num{0.3}n}_\text{Svist} = \num{0.23}n
      \end{equation*}
      funkčních raket. Pravděpodobnost zakoupení funkční rakety je tedy
      \begin{equation*}
        p(A) = \dfrac{m}{n} = \num{0.23}.
      \end{equation*}
      Výsledek lze snadno zobecnit. Označíme-li \(p_1\) pravděpodobnost zmetku ve firmě Hvizd a
      \(p_2\) pravděpodobnost zmetku ve firmě Svist, je pravděpodobnost funkčního výrobku ve Hvizdu
      \((1 - p_1)\) a ve Svistu \((1 - p_2)\). Označme \(q\) podíl Hvizdu na celkové produkci. Podíl
      Svistu je pak \((1 - q)=\num{0.3}\). Pravděpodobnost koupě funkčního výrobku je
      \begin{align*}
        p(A) &= (1 - p_1)q + (1 - p_2)(1 - q)    \\
             &= (1 - p_2) + q(p_2 - p_1).
      \end{align*}
      
      V úloze b) se již jedná o \textbf{podmíněnou pravděpodobnost}. Koupíme v obchodě raketu,
      podíváme se na obal a zjistíme, že byla vyrobena ve Hvizdu. S touto dodatečnou informací
      chceme zjistit pravděpodobnost, že až raketu rozbalíme a odpálíme, bude skutečně fungovat.
      Označme jako jev \(B\) „Raketa byla vyrobena ve Hvizdu.“ Naším úkolem je tedy zjistit
      pravděpodobnost jevu \(A\) (raketa bude funkční) za podmínky, že nastal jev \(B\) (byla
      vyrobena ve Hvizdu). Tuto pravděpodobnost značíme \(p_B(A)\). Víme, že na trhu je \(qn =
      \num{0.7}n\) raket vyrobených ve Hvizdu. To představuje pro náš další výpočet počet případů
      možných. Pouze \((1 - p_1) = \num{0.17}\) z nich je však funkčních, počet případů příznivých
      je tedy \((1 - p_1)qn = \num{0.17}\cdot\num{0.7}n = \num{0.119}n\). Hledaná pravděpodobnost je
      \begin{equation*}
        p_B(A) = \dfrac{(1 - p_1)qn}{qn} = 1 - p_1 = \num{0.17}
      \end{equation*}
      Všimněme si výpočtu podrobněji. \((1 - p_1)q\) představuje pravděpodobnost jevu \((A\text{ a
      }B)\), že náhodně zakoupená raketa bude funkční a zároveň bude vyrobena ve Hvizdu. Skute čně,
      na trhu je v dané chvíli \(n\) raket, z nich \(qn\) bylo vyrobeno ve Hvizdu a z těchto \(qn\)
      výrobků Hvizdu je \((1 - p_1)qn\) funkčních. Proto 
      \begin{equation*}
        p(A\text{ a }B) = \dfrac{(1 - p_1)qn}{n} = (1 - p_1)q = \num{0.119}
      \end{equation*}
      (Proč pravděpodobnost není součinem \(p(A)p(B)\)? Nedivme se, jevy \(A\) a \(B\) nejsou totiž
      nezávislé!). Vidíme, že platí
      \begin{equation*}
        \boxed{p(A\text{ a }B) = p(B)p_B(A)}\,,
      \end{equation*}
      neboť \(q = p(B)\). Získáváme tedy vztah pro výpočet podmíněné pravděpodobnosti:
      \begin{mdframed}[style=highlight]
        \begin{equation}\label{mai:eq057}
          p_B(A) = \dfrac{p(A\text{ a }B)}{p(B)} \quad\text{obdobně}\quad
          p_A(B) = \dfrac{p(A\text{ a }B)}{p(A)}.
        \end{equation}
      \end{mdframed}
    
      Zbývá vyřešit soutěžní úkol c). Koupili jsme raketu a odpálili ji. Fungovala. Jaká je
      pravděpodobnost, že když se nyní podíváme na obal, zjistíme, že šlo o výrobek Svistu? Můžeme
      již využít vztahu (\ref{mai:eq057}). Máme totiž počítat pravděpodobnost jevu \(B\) za
      podmínky, že nastal jev \(A\). (Poznamenejme, že je-li jev \(B\) definován jako „Náhodně
      koupená raketa pochází z Hvizdu.“, je jev „Náhodně koupená raketa pochází ze Svistu.“ jevem
      opačným k \(B\) tj. \(p(\overline{B})=1-q\).) Platí
      \begin{align*}
        p_A(\overline{B})&= \dfrac{p(A\text{ a }\overline{B})}{p(A)}
                          = \dfrac{p_{\overline{B}}(A)p(\overline{B})}{p(A)}                \\             
                         &= \dfrac{(1 - p_2)(1 - q)}{(1 - p_2) + q(p_2 - p_1)}              \\
                         &= \dfrac{\num{0.37}\cdot\num{0.3}}{\num{0.37} + \num{0.7}\cdot(\num{-0.2})}
                          \simeq\num{0.483}
      \end{align*}
      Kdo chce, může snadno dospět k výsledku přímou úvahou: Na trhu je \(n\) raket, z nich
      funkčních je \(\num{0.17}\cdot\num{0.7}n + \num{0.37} - \num{0.3}n = \num{0.23}n\). Tato
      hodnota je při výpočtu \(p(A\text{ a }B)\)  počtem případů možných. Počet funkčních raket
      vyrobených Svistem je \(\num{0.37}\cdot\num{0.3}n = \num{0.111}n\). To je počet případů
      příznivých. Hledaná pravděpodobnost je tedy podílem
      \begin{equation*}
        p(A\text{ a }\overline{B}) = \dfrac{\num{0.111}n}{\num{0.23}n}\simeq\num{0.483}.
      \end{equation*}
  \end{example}
\end{mdframed}